\section{一个唯一性定理}

在本小节中，对任何严格正的整数 m、区间 \([1,m]\) 中的任何整数 i 以及任何序列 \((v_{1},\ldots,v_{m})\)（在 \(V^{m}\) 中），我们用 \((v_{1},\ldots,\widehat{v_{i}},\ldots,v_{m})\) 表示序列 \((v_{1},\ldots,v_{i-1},v_{i+1},\ldots,v_{m})\)（在 \(V^{m-1}\) 中）。

命题 1. --- 设 W 是 V 中的一个超平面，e 是 V---W 中的一个向量。设 \(\varphi\) 是 \(\Omega(\mathrm{W})\) 的一个元素。存在唯一的元素 \(\omega\) 属于 \(\Omega(\mathrm{V})\)，使得我们有

\[
\omega (w _ {1}, \ldots , w _ {n - 1}, e) = \varphi (w _ {1}, \ldots , w _ {n - 1})\tag{7}
\]

对 W 的每个基 \((w_{1},\ldots ,w_{n - 1})\) 成立。

\begin{enumerate}
\def\labelenumi{\Alph{enumi})}
\tightlist
\item
  \(\omega\) 的唯一性：设 \(\omega\) 是 \(\Omega(\mathrm{V})\) 的一个元素。设 \((v_{1},\ldots,v_{n})\) 是 V 的一个基，\(\mu_{1},\ldots,\mu_{n}\) 是 e 在这个基中的坐标。用 I 表示从 1 到 n（含端点）中使 \(\mu_{i}\neq0\) 的整数所成的集合；它非空。设 i 是 I 的一个元素。序列 \((v_{1},\ldots,\widehat{v_{i}},\ldots,v_{n},e)\) 是 V 的一个基，且由于 \(e-v_{i}\mu_{i}\) 是序列 \((v_{1},\ldots,\widehat{v_{i}},\ldots,v_{n})\) 的一个线性组合，公式 (6) 蕴含
\end{enumerate}

\[
\omega (v _ {1}, \ldots , \widehat {v _ {i}}, \ldots , v _ {n}, v _ {i} \mu_ {i}) = \omega (v _ {1}, \ldots , \widehat {v _ {i}}, \ldots , v _ {n}, e).\tag{8}
\]

由公式 (1) 和 (3) 可知，这个等式的左边等于 \(\pi(-1)^{n-i}\pi(\mu_i)\omega(v_1,\ldots,v_n)\)。用 p 表示 V 的以 W 为像、以 eD 为核的投影。向量 \(v_j - p(v_j)\) 都与 e 成比例，反复应用公式 (6) 表明公式 (8) 的右边等于 \(\omega(p(v_1),\ldots,\widehat{p(v_i)},\ldots,p(v_n),e)\)。

由此得出，若 \(\omega\) 满足关系式 (7)，则我们有

\[
\omega (v _ {1}, \dots , v _ {n}) = \pi (- 1) ^ {n - i} \pi (\mu_ {i}) ^ {- 1} \varphi (p (v _ {1}), \dots , \widehat {p (v _ {i})}, \dots , p (v _ {n}))\tag{9}
\]

对 I 中的每个 i 成立，这就证明了 \(\omega\) 的唯一性。

\begin{enumerate}
\def\labelenumi{\Alph{enumi})}
\setcounter{enumi}{1}
\tightlist
\item
  \(\omega\) 的构造：设 \((v_{1},\ldots,v_{n})\) 是 V 的一个基；我们像上面那样定义 \(\mu_{1},\ldots,\mu_{n}\) 与 I。对 I 中的每个 \(i\)，假设 \(\mu_{i}\neq0\) 蕴含序列 \((v_{1},\ldots,\widehat{v_{i}},\ldots,v_{n},e)\) 生成空间 V，因此序列 \((p(v_{1}),\ldots,\widehat{p(v_{i})},\ldots,p(v_{n}))\) 生成空间 W = p(V)。换言之，后一序列是 W 的一个基，且
\end{enumerate}

\[
t _ {i} = \pi (- 1) ^ {n - i} \pi (\mu_ {i}) ^ {- 1} \varphi (p (v _ {1}), \ldots , \widehat {p (v _ {i})}, \ldots , p (v _ {n}))\tag{10}
\]

定义了 \(\mathrm{D}_{\mathrm{ab}}^{*}\) 的一个元素。

设 i 与 j 是 I 的两个满足 i \textless{} j 的元素；让我们证明等式 \(t_{i} = t_{j}\)。按定义，向量 \(v_{i}\mu_{i} + v_{j}\mu_{j} - e\) 是诸向量 \(p(v_{k})\)（k 异于 i 与 j）的一个线性组合。因此，\(p(v_{i})\mu_{i} + p(v_{j})\mu_{j}\) 是诸向量 \(p(v_{k})\)（k 异于 i 与 j）的一个线性组合。于是由公式 (6)，我们有

\[
\begin{array}{c} \varphi (p (v _ {1}), \ldots , \widehat {p (v _ {i})}, \ldots , \widehat {p (v _ {j})}, \ldots , p (v _ {n}), p (v _ {i}) \mu_ {i}) \\ = \varphi (p (v _ {1}), \ldots , \widehat {p (v _ {i})}, \ldots , \widehat {p (v _ {j})}, \ldots , p (v _ {n}), - p (v _ {j}) \mu_ {j}). \end{array}
\]

应用公式 (1) 与 (3)，我们推出

\[
\begin{array}{c} \varphi (p (v _ {1}), \ldots , \widehat {p (v _ {j})}, \ldots , p (v _ {n})) \pi (\mu_ {i}) \pi (- 1) ^ {n - i - 1} \\ = \varphi (p (v _ {1}), \ldots , \widehat {p (v _ {i})}, \ldots , p (v _ {n})) \pi (- \mu_ {j}) \pi (- 1) ^ {n - j}, \end{array}
\]

从而 \(t_i = t_j\)。

证明了这一点之后，我们把 \(\omega(v_1, \ldots, v_n)\) 定义为诸 \(t_i\)（\(i\) 属于 I）的公共值。这样我们就构造了一个映射 \(\omega\)（从 B(V) 到 D \(_{\text{ab}}^*\)），它满足关系式 (9)。设 \((w_1, \ldots, w_{n-1})\) 是 W 的一个基；若置 \(v_i = w_i\)（对 \(1 \leqslant i \leqslant n-1\)）且置 \(v_n = e\)，则我们有 I = \{n\} 与 \(\mu_n = 1\)。我们还有 \(p(v_i) = w_i\)（对 \(1 \leqslant i \leqslant n-1\)），而公式 \(\omega(w_1, \ldots, w_{n-1}, e) = \varphi(w_1, \ldots, w_{n-1})\) 是 (9) 的一个特殊情形。

现在我们必须证明 \(\omega\) 属于 \(\Omega (\mathrm{V})\)。

\begin{enumerate}
\def\labelenumi{\Alph{enumi})}
\setcounter{enumi}{2}
\tightlist
\item
  公式 (1) 的证明：设 \(\lambda_1, \ldots, \lambda_n\) 是 \(D^*\) 的元素，\((v_1, \ldots, v_n)\) 是 \(V\) 的一个基。我们像上面那样定义 \(\mu_1, \ldots, \mu_n\) 与 \(I\)。取定一个 \(i\) 属于 \(I\)。由于我们有 \(e = \sum_{j=1}^{n} (v_j \lambda_j)(\lambda_j^{-1} \mu_j)\)，公式 (9) 蕴含
\end{enumerate}

\[
\begin{array}{l l} (1 1) & \omega (v _ {1} \lambda_ {1}, \ldots , v _ {n} \lambda_ {n}) \\ & \qquad = \pi (- 1) ^ {n - i} \pi (\lambda_ {i} ^ {- 1} \mu_ {i}) ^ {- 1} \varphi (p (v _ {1}) \lambda_ {1}, \ldots , \widehat {p (v _ {i}) \lambda_ {i}}, \ldots , p (v _ {n}) \lambda_ {n}). \end{array}
\]

但我们有 \(\pi (\lambda_i^{-1}\mu_i)^{-1} = \pi (\mu_i)^{-1}\pi (\lambda_i)\)，且

\[
\begin{array}{c} \varphi (p (v _ {1}) \lambda_ {1}, \ldots , \widehat {p (v _ {i}) \lambda_ {i}}, \ldots , p (v _ {n}) \lambda_ {n}) \\ = \varphi (p (v _ {1}), \ldots , \widehat {p (v _ {i})}, \ldots , p (v _ {n})) \pi (\lambda_ {1} \ldots \widehat {\lambda_ {i}} \ldots \lambda_ {n}). \end{array}
\]

比较公式 (9) 与 (11) 就得到所要的关系式

\[
\omega (v _ {1} \lambda_ {1}, \ldots , v _ {n} \lambda_ {n}) = \omega (v _ {1}, \ldots , v _ {n}) \pi (\lambda_ {1} \ldots \lambda_ {n}).
\]

\begin{enumerate}
\def\labelenumi{\Alph{enumi})}
\setcounter{enumi}{3}
\tightlist
\item
  公式 (2) 的证明：设 \((v_{1},\ldots,v_{n})\) 是 V 的一个基，i、j 是区间 {[}1,n{]} 中两个不同的整数；我们像先前那样定义 \(\mu_{1},\ldots,\mu_{n}\)。考虑 V 的基 \((v_{1}^{\prime},\ldots,v_{n}^{\prime})\)，它由 \(v_{i}^{\prime}=v_{i}+v_{j}\) 与 \(v_{k}^{\prime}=v_{k}\)（\(k\neq i\)）定义，并引入 e 在这个基中的坐标 \(\mu_{1}^{\prime},\ldots,\mu_{n}^{\prime}\)；它们满足 \(\mu_{j}^{\prime}=\mu_{j}-\mu_{i}\) 与 \(\mu_{k}^{\prime}=\mu_{k}\)（\(k\neq j\)）。我们置 \(t=\omega(v_{1},\ldots,v_{n})\) 与 \(t^{\prime}=\omega(v_{1}^{\prime},\ldots,v_{n}^{\prime})\)。我们必须证明 t 与 \(t^{\prime}\) 相等。
\end{enumerate}

首先注意，按 \(\omega\) 的定义，我们有

\begin{enumerate}
\def\labelenumi{(\arabic{enumi})}
\setcounter{enumi}{11}
\tightlist
\item
\end{enumerate}

\[
t = \pi (- 1) ^ {n - k} \pi (\mu_ {k}) ^ {- 1} \varphi (p (v _ {1}), \ldots , \widehat {p (v _ {k})}, \ldots , p (v _ {n})),\tag{13}
\]

\[
t ^ {\prime} = \pi (- 1) ^ {n - k} \pi (\mu_ {k} ^ {\prime}) ^ {- 1} \varphi (p (v _ {1} ^ {\prime}), \ldots , \widehat {p (v _ {k} ^ {\prime})}, \ldots , p (v _ {n} ^ {\prime}))
\]

对每个 \(k\)（使 \(\mu_{k}\) 与 \(\mu_k^\prime\) 都非零）成立。

\begin{enumerate}
\def\labelenumi{\alph{enumi})}
\item
  若 \(\mu_i \neq 0\)，则序列 \((p(v_1), \ldots, \widehat{p(v_k)}, \ldots, p(v_n))\) 与 \((p(v_1'), \ldots, \widehat{p(v_k'}, \ldots, p(v_n'))\) 相等，且我们有 \(\mu_i = \mu_i'\)，从而 \(t = t'\)。
\item
  若存在一个指标 \(k\)（异于 \(i\) 与 \(j\)）使 \(\mu_k \neq 0\)，则我们有 \(p(v_l') = p(v_l)\)（对 \(l \neq i\)），且 \(p(v_i') = p(v_i) + p(v_j)\)。元素 \(p(v_i)\)
\end{enumerate}

与 \(p(v_{j})\) 都属于序列 \((p(v_{1}),\ldots ,\widehat{p(v_{k})},\ldots ,p(v_{n}))\)。由于 \(\varphi\) 属于 \(\Omega (\mathrm{W})\)，VIII, 第 447 页的公式 (2) 适用；它给出

\[
\varphi (p (v _ {1}), \dots , \widehat {p (v _ {k})}, \dots , p (v _ {n})) = \varphi (p (v _ {1} ^ {\prime}), \dots , \widehat {p (v _ {k} ^ {\prime})}, \dots , p (v _ {n} ^ {\prime})).
\]

但我们还有 \(\mu_{k} = \mu_{k}^{\prime}\)，从而 \(t = t'\)。

\begin{enumerate}
\def\labelenumi{\alph{enumi})}
\setcounter{enumi}{2}
\tightlist
\item
  只剩下考察指标 \(k\)（使 \(\mu_k \neq 0\) 的唯一指标）是 \(j\) 的情形。这时我们有 \(\mu_j' = \mu_j\)、\(e = v_j\mu_j\) 与 \(p(v_j) = 0\)。在公式 (12) 与 (13) 中取 \(k = j\)。由于序列 \((p(v_1), \ldots, \widehat{p(v_j)}, \ldots, p(v_n))\) 与 \((p(v_1'), \ldots, \widehat{p(v_j'), \ldots, p(v_n'))}\) 相等，我们有 \(t = t'\)。
\end{enumerate}

推论 1. --- 设 \((e_1, \ldots, e_n)\) 是 V 的一个基，t 是 \(\mathrm{D}_{\mathrm{ab}}^*\) 的一个元素。存在唯一的元素 \(\omega\) 属于 \(\Omega(\mathrm{V})\)，使得 \(\omega(e_1, \ldots, e_n) = t\)。

我们对 n 用归纳法证明该推论。若 n = 0，则 V 的唯一的基是空基，故断言为真。现在设 \(n \geqslant 1\)；用 W 表示由 \(e_{1}, \ldots, e_{n-1}\) 生成的 V 的子空间。由命题 1，存在一个双射 \(\Lambda\)（从 \(\Omega(\mathrm{V})\) 到 \(\Omega(\mathrm{W})\)），满足

\[
(\Lambda (\omega)) (w _ {1}, \ldots , w _ {n - 1}) = \omega (w _ {1}, \ldots , w _ {n - 1}, e _ {n})
\]

对 W 的每个基 \((w_{1},\ldots,w_{n-1})\) 与每个 \(\omega\)（属于 \(\Omega(\mathrm{V})\)）成立。由归纳假设，存在唯一的元素 \(\varphi\) 属于 \(\Omega(\mathrm{W})\)，使得 \(\varphi(e_{1},\ldots,e_{n-1})=t\)。关系式 \(\omega(e_{1},\ldots,e_{n})=t\)（\(\omega\) 属于 \(\Omega(\mathrm{V})\)）等价于 \(\Lambda(\omega)=\varphi\)。由此即得推论 1。

推论 2. --- 设 \(\omega\) 与 \(\omega'\) 是 \(\Omega(V)\) 的两个元素。存在唯一的元素 \(t\) 属于 \(\mathrm{D}_{\mathrm{ab}}^{*}\)，使得 \(\omega' = t\omega\)。

注. --- 假设体 D 是交换的。按定义，B(V) 是 \(V^{n}\) 的一个子集。显然，非零交错 n 重线性形式 \(f: V^{n} \to D\) 在 B(V) 上的限制属于 \(\Omega(V)\)。此外，若 \((e_{1}, \ldots, e_{n})\) 是 V 的一个基且 t 是 D 的一个非零元素，则存在唯一的交错 n 重线性形式 f，使得 \(f(e_{1}, \ldots, e_{n}) = t\)（III, §7, No.~4, 第 511 页与 III, §7, No.~8, 第 518 页）。由推论 1，集合 \(\Omega(V)\) 由诸非零交错 n 重线性形式在 B(V) 上的限制组成。

